\subsection{Колоночные файловые форматы хранения данных}\label{subsec:file-formats}

Файловый формат задаёт, как табличные данные раскладываются в байты внутри файла.
Принципиальное различие --- построчное (row-oriented) и поколоночное (columnar)
хранение. При построчном хранении значения одной строки лежат рядом, что выгодно для
точечных операций по ключу (прочитать или записать целую запись) и для потоковой
записи (строка добавляется в конец без переупорядочивания). При поколоночном
хранении рядом лежат значения одного столбца: это даёт высокую степень сжатия
(однотипные данные сжимаются лучше), позволяет читать только запрашиваемые столбцы
(column pruning) и эффективно применять векторизованные операции к плотным массивам
значений. Аналитические нагрузки (OLAP) --- сканирование, фильтрация и агрегация по
немногим столбцам широких таблиц --- естественно ложатся на колоночное хранение,
поэтому далее рассматриваются именно колоночные форматы~\cite{stonebraker2005cstore, abadi2013column}.

Apache Parquet~\cite{parquet-format} --- де-факто стандарт колоночного хранения в экосистеме
больших данных. Файл делится на группы строк (row groups, типично десятки или сотни
мегабайт); внутри группы строк каждый столбец хранится отдельным фрагментом (column
chunk), который, в свою очередь, разбивается на страницы (pages, типично около
\SI{1}{\mega\byte}). Страница --- единица сжатия и кодирования. Parquet применяет
словарное кодирование (dictionary encoding) для столбцов с небольшим числом
различных значений, кодирование длин серий (run-length encoding) и упаковку в
минимально необходимое число бит (bit-packing) для целочисленных и словарных
индексов, дельта-кодирование для возрастающих последовательностей, а поверх ---
блочное сжатие (Snappy, ZSTD, GZIP). Для пропуска данных Parquet хранит статистики
(минимум, максимум, число null) на уровне группы строк и страницы, а также
опционально --- фильтры Блума~\cite{bloom1970space}. Чтение в JVM-окружении выполняется библиотекой
\texttt{parquet-mr}; в Apache Paimon используется её затенённая (shaded) сборка.

Apache ORC~\cite{orc-spec} (Optimized Row Columnar) близок по идеологии к Parquet, но
исторически развивался в экосистеме Apache Hive. Файл делится на полосы (stripes,
типично около \SI{64}{\mega\byte}); каждая полоса содержит данные столбцов, индексные
данные и нижний колонтитул. ORC применяет «лёгкое» кодирование (RLE для целых,
словарное для строк) и «тяжёлое» блочное сжатие (ZLIB, Snappy, ZSTD); ведёт
встроенные индексы со статистиками на уровне полосы и группы строк (по \num{10000}
строк), что делает его особенно эффективным на точечных и узкодиапазонных запросах.
Реализация --- нативная для JVM библиотека \texttt{orc-core}.

Vortex~\cite{vortex-repo, vortex-docs} --- современный колоночный формат, разрабатываемый на языке Rust с
прицелом на нативное векторизованное выполнение. Ключевые особенности. Во-первых,
система кодирования: целочисленные данные кодируются по схеме FastLanes~\cite{afroozeh2023fastlanes} ---
бит-упаковка, ориентированная на SIMD-декодирование (значения раскладываются так,
чтобы распаковка выполнялась векторными инструкциями процессора без ветвлений);
числа с плавающей точкой --- кодеком ALP~\cite{afroozeh2024alp} (Adaptive Lossless floating-Point), который
без потерь представляет вещественные значения как целые числа с общим масштабом, а
исключения хранит отдельно; строки --- кодеком FSST~\cite{boncz2020fsst} (Fast Static Symbol Table),
строящим словарь коротких подстрок (символов) и заменяющим повторяющиеся фрагменты
ссылками на словарь, что особенно эффективно на строках с общими префиксами (URL,
идентификаторы, пути). Кодеки в Vortex --- расширяемые и каскадируемые: к закодированному
массиву можно применить ещё один кодек. Во-вторых, физическая модель данных Vortex
совместима с Apache Arrow~\cite{arrow-site}, что позволяет передавать декодированные пакеты в движок
без копирования через интерфейс Arrow C Data Interface. В-третьих, Vortex
поддерживает проталкивание предикатов (predicate pushdown) на уровне нативного
декодера: условие фильтрации преобразуется в нативное выражение, и блоки данных,
заведомо не содержащие подходящих строк, не декодируются. В-четвёртых, метаданные
файла (схема, статистики, расположение блоков) выделены в отдельный «футер»,
переносимый и читаемый независимо от данных. Доступ к нативной библиотеке Vortex из
JVM осуществляется через Java Native Interface (JNI).

Lance~\cite{lance-docs} --- колоночный формат, также реализованный на Rust, но
ориентированный на нагрузки машинного обучения: быстрый произвольный доступ к
строкам (random access), хранение векторных представлений (embeddings) и встроенные
векторные индексы. Lance оптимизирует чтение отдельных строк и небольших выборок,
что важно при обучении моделей и поиске ближайших соседей, но менее существенно для
классической OLAP-аналитики со сканированием. Поддержка Lance в Apache Paimon на
момент работы существует, однако требует доработки для корректного слияния версий в
таблицах с первичным ключом. В предварительном прогоне near-real-time-сценария Lance
измерялся, но оказался стабильно худшим форматом по времени чтения (в \num{2}--\num{4}
раза медленнее ORC и Parquet, заметнее всего --- на точечных запросах, где его
ориентация на произвольный доступ не реализуется из-за слияния версий при чтении),
поэтому в итоговое сравнение второго сценария он не вошёл. На пакетной аналитике
TPC-DS он также оказался самым медленным
из четырёх форматов (в среднем в \num{1{,}5}--\num{2} раза медленнее по времени
запросов и в \num{8} раз больше по объёму хранилища), что подтверждает его
ориентацию на иную нагрузку --- произвольный доступ к строкам и векторный поиск,
а не сканирование.

Apache Avro~\cite{avro-spec} в сравнение не включён принципиально: это строчный
формат. Avro хранит записи последовательно, со схемой в заголовке файла; он удобен
для потоковой передачи и журналирования (дёшев на запись, поддерживает эволюцию
схемы), но не даёт ни поколоночного сжатия, ни чтения отдельных столбцов, ни
статистик для пропуска данных. Сравнивать его с колоночными форматами на
аналитической нагрузке некорректно. Avro упоминается в работе лишь в связи с
гипотезой о его использовании для самых свежих, ещё не уплотнённых файлов
LSM-дерева; эта гипотеза рассмотрена и отклонена в практической части.

Сводная характеристика рассматриваемых форматов приведена в
таблице~\ref{tab:formats-overview}.

\begin{xltabular}{\textwidth}{|l|X|X|}
    \caption{Сравнительная характеристика колоночных файловых форматов}\label{tab:formats-overview} \\ \hline
    \textbf{Формат} & \textbf{Реализация чтения} & \textbf{Отличительные черты} \\ \hline
    \endfirsthead
    \caption*{Продолжение таблицы~\ref{tab:formats-overview}} \\ \hline
    \textbf{Формат} & \textbf{Реализация чтения} & \textbf{Отличительные черты} \\ \hline
    \endhead
    \hline
    \endfoot
    \hline
    \endlastfoot
    Parquet & JVM (\texttt{parquet-mr}, затенённая сборка в Paimon) &
        группы строк и страницы; словарное / RLE / bit-packing кодирование; блочное сжатие; статистики и фильтры Блума; зрелый универсальный формат \\ \hline
    ORC & JVM (\texttt{orc-core}) &
        полосы (stripes); лёгкое и тяжёлое сжатие; встроенные индексы по \num{10000} строк; преимущество на точечных и узкодиапазонных запросах \\ \hline
    Vortex & нативная (Rust) через JNI &
        FastLanes (SIMD-bit-packing), ALP (числа с плавающей точкой), FSST (строки по схожести префиксов); каскадируемые кодеки; совместимость с Arrow; нативное проталкивание предикатов; переносимый футер метаданных \\ \hline
    Lance & нативная (Rust) через JNI &
        быстрый произвольный доступ к строкам; векторные индексы; ориентация на нагрузки машинного обучения \\ \hline
    Avro & JVM &
        строчный формат; дёшев на запись; нет поколоночного сжатия, column pruning и статистик; в сравнение не включён \\ \hline
\end{xltabular}
